\chapter[Introduction]{Introduction}
\label{Chap:Intro}
\counterwithout{table}{chapter}
\counterwithout{figure}{chapter}
\makeatletter
\setlength{\@fptop}{0pt}
\makeatother

Les galaxies sont des systèmes liés gravitationnellement, constitués principalement d'étoiles, de gaz et de poussière. Leur évolution dépend étroitement des interactions entre ces composantes. La formation stellaire joue un rôle fondamental dans cette évolution : elle détermine la croissance des populations stellaires et modifie les propriétés physiques des galaxies au cours du temps cosmique. Elle est alimentée par le gaz du milieu interstellaire (MIS), qui occupe l'espace entre les étoiles. Les étoiles naissent dans des nuages moléculaires froids et denses, composés principalement d'hydrogène moléculaire (H$_2$). Cette phase moléculaire se forme à partir du gaz atomique diffus par des transitions thermodynamiques et chimiques qui ne sont pas encore entièrement comprises \citep{Kennicutt1998, Klessen2016}. Le rythme auquel se forment les étoiles, généralement exprimé par le taux de formation stellaire (SFR, pour \textit{star formation rate}), est donc étroitement lié aux propriétés physiques et à l'état thermodynamique des différentes phases du gaz interstellaire \citep{Fbigiel2008, Leroy2008}. Ce mémoire étudie les liens entre la formation stellaire, le gaz moléculaire et les phases de l'\HI.\\\indent
Le passage de l'hydrogène atomique à l'hydrogène moléculaire constitue une étape importante entre le gaz diffus et la formation des étoiles. L'hydrogène moléculaire se forme essentiellement par des réactions catalytiques à la surface des grains de poussière interstellaire, sa formation en phase gazeuse étant peu efficace dans les conditions interstellaires habituelles \citep{Gould1963, Hollenbach1971}. Une fois formées, les molécules de H$_2$ peuvent être photodissociées par le rayonnement ultraviolet (UV) dans les bandes de Lyman--Werner. Leur survie dépend donc de leur protection contre ce rayonnement, assurée par l'absorption par la poussière et, lorsque la densité de colonne est suffisamment élevée, par l'auto-écrantage du H$_2$ \citep{Draine1996, Krumholz2009, Sternberg2014}.\\\indent
Les modèles théoriques prévoient que cette protection devient efficace lorsqu'une densité de colonne suffisante s'est accumulée \citep{Krumholz2009, Sternberg2014}. La phase froide et dense de l'\HI, appelée milieu neutre froid (CNM, pour \textit{cold neutral medium}), possède une densité relativement élevée ($\sim$10--100 cm$^{-3}$). Elle peut ainsi atteindre plus facilement la densité de colonne nécessaire à la protection des molécules que la phase tiède, appelée milieu neutre tiède (WNM, pour \textit{warm neutral medium}), dont la densité est de l'ordre de $0.1$--1 cm$^{-3}$ \citep{Krumholz2009, Sternberg2014, Syed2020}. De plus, le taux de formation du H$_2$ sur les grains dépend du produit des densités en nombre de l'hydrogène atomique et des grains de poussière : les régions de CNM les plus denses convertissent donc plus efficacement le gaz atomique en gaz moléculaire \citep{Hollenbach1979, Wakelam2017}. Le CNM constitue ainsi à la fois un environnement propice à la protection des molécules et un réservoir de gaz atomique à partir duquel les nuages moléculaires peuvent se former \citep{Wolfire2003, Klessen2016}.\\\indent
Malgré l'importance théorique du CNM dans la formation des nuages moléculaires, les études observationnelles des liens entre la formation stellaire et le MIS se sont surtout intéressées à l'\HI\ total ou au gaz moléculaire tracé par l'émission du monoxyde de carbone (CO) \citep{Kennicutt1998, Fbigiel2008, Leroy2008}. Cela tient notamment à la difficulté de séparer le CNM du WNM par les observations. Les deux phases contribuent à la même raie d'émission à 21 cm. Les mesures d'absorption devant des sources de continuum d'arrière-plan sondent préférentiellement le gaz froid, mais ne fournissent qu'un échantillonnage spatial limité. La décomposition spectrale offre une autre possibilité, dont les développements récents permettent désormais une application systématique aux galaxies extérieures \citep{Heiles2003a, Marchal2019, Syed2020}.\\\indent
Comprendre les liens entre les phases du MIS et la formation stellaire est essentiel pour décrire l'évolution des galaxies. Les galaxies proches constituent des laboratoires privilégiés, car elles peuvent être observées dans leur ensemble, contrairement à la Voie lactée, où notre position à l'intérieur du disque limite notre vue \citep{Leroy2008}. Établir des relations quantitatives entre le CNM et le gaz moléculaire fournirait également un diagnostic utile dans les environnements où les observations du CO sont absentes ou peu fiables. C'est particulièrement le cas des galaxies naines de faible métallicité : la moindre protection par la poussière y favorise la photodissociation du CO, alors que le H$_2$ peut subsister grâce à son auto-écrantage \citep{Bolatto2013}.\\\indent
Les relations en loi de puissance entre la densité de surface du gaz et celle du taux de formation stellaire, généralement appelées lois de Kennicutt--Schmidt, restent parmi les outils empiriques les plus utilisés pour relier le gaz à la formation des étoiles \citep{Schmidt1959, Kennicutt1998}. Le rôle de chaque phase du MIS demeure cependant mal contraint par les observations. Dans leur forme la plus courante, ces relations regroupent le gaz neutre en un seul réservoir, en additionnant l'\HI\ et le H$_2$. Il devient alors difficile d'identifier la phase qui gouverne la relation. Les observations résolues de THINGS montrent que la densité de surface du taux de formation stellaire varie à peu près linéairement avec $\Sigma_{\mathrm{H_2}}$, tandis que l'\HI\ sature vers $9$--$10\ \mathrm{M_\odot\,pc^{-2}}$ et ne présente plus qu'une faible corrélation avec la formation stellaire au-delà de ce seuil \citep{Fbigiel2008}. Des simulations du milieu neutre froid dans les disques spiraux montrent, quant à elles, une relation entre la densité du taux de formation stellaire et la densité de colonne du CNM, semblable à la loi de Kennicutt--Schmidt du gaz total. Le CNM y trace mieux le gaz associé à la formation stellaire que l'ensemble du gaz atomique neutre \citep{Smith2023}. Cette étude vise à caractériser les relations entre les densités de surface du gaz et du taux de formation stellaire dans un échantillon de galaxies proches, en examinant plus particulièrement si la composante cinématiquement froide de l'\HI\ peut tracer les conditions qui précèdent la formation des nuages moléculaires et des étoiles. Nous appliquons des méthodes de décomposition spectrale à des cubes de données \HI\ de haute résolution afin de quantifier ses liens avec le gaz moléculaire et l'activité de formation stellaire. Cette composante cinématique sert de traceur indirect du CNM défini thermiquement. Nous cherchons ainsi à tester la prédiction selon laquelle l'abondance du CNM régule l'efficacité de la transition \HI--H$_2$ et de la formation stellaire qui s'ensuit.

\section{Le milieu interstellaire}
\subsection{Vue d'ensemble}

L'espace entre les étoiles n'est pas vide : il contient un mélange de gaz, de poussière, de particules chargées relativistes, ou rayons cosmiques, et de champs électromagnétiques. Cet ensemble constitue le milieu interstellaire \citep{ISF2011, Ferriere_2001}.\\\indent
Le gaz représente la majeure partie de la masse du MIS et se concentre généralement dans des structures nuageuses de tailles et de morphologies variées \citep{ISF2011}. La figure~\ref{fig:monkeyhead} montre, par exemple, un nuage d'hydrogène moléculaire. L'hydrogène est l'élément le plus abondant dans les galaxies ; il s'y trouve sous forme atomique (\HI), moléculaire (H$_2$) ou ionisée (H\textsc{ii}). Dans la Voie lactée, le MIS comprend en masse environ 70\% d'hydrogène, 28\% d'hélium et 2\% d'éléments plus lourds \citep{Ferriere_2001}. En astronomie, tous les éléments plus lourds que l'hélium, dont le carbone, l'azote et l'oxygène, sont désignés collectivement par le terme de « métaux ».\\\indent
La poussière interstellaire se compose de petits grains solides, généralement de taille inférieure au micromètre, constitués principalement de silicates, de matériaux carbonés et de composés contenant du fer \citep{Zhukovska_2008}. Certains nuages moléculaires apparaissent sombres sur le fond stellaire, car leur poussière absorbe et diffuse la lumière visible des étoiles situées derrière eux. La nébuleuse du Sac à charbon (figure~\ref{fig:coalsack}) en fournit un exemple caractéristique \citep{ISF2011}.\\\indent
Les propriétés et la dynamique du MIS présentent un intérêt majeur, car les cycles de vie des étoiles et de la matière interstellaire sont étroitement liés. Les étoiles naissent de l'effondrement gravitationnel des nuages de gaz. Durant la séquence principale, elles transforment l'hydrogène en hélium dans leur cœur. À la fin de leur évolution, lorsque leur combustible en hydrogène est presque épuisé, elles éjectent leurs couches externes sous forme de nébuleuses planétaires ou, pour les étoiles les plus massives, explosent en supernovæ. Dans les deux cas, la matière enrichie en éléments lourds retourne au MIS et peut participer à la formation de nouvelles générations d'étoiles \citep{ISF2011}. La dynamique du MIS détermine ainsi où et quand les étoiles se forment, tandis que les propriétés des étoiles et de leurs systèmes planétaires dépendent des conditions physiques et chimiques de leur environnement natal.

\subsection{Les phases du gaz interstellaire}

Le gaz interstellaire est loin d'être homogène : ses propriétés physiques varient fortement selon la région considérée dans une galaxie. Aucun classement unique ne permet donc de décrire toutes ses caractéristiques. Plusieurs cadres ont été proposés pour en distinguer les phases. \citet{Klessen2016} en décrivent cinq, résumées dans le tableau ci-dessous.

\begin{table}[htbp]
\centering
\caption{\textbf{Description et propriétés physiques typiques des phases du MIS.} Valeurs issues de \citet{Klessen2016} et de \citet{Ferriere_2001}.}
\label{tab:ISMcomponents}
\small
\setlength{\tabcolsep}{6pt}
\renewcommand{\arraystretch}{1.15}
\begin{tabularx}{\linewidth}{L{3.2cm} L{6.3cm} c c c}
\toprule
Phase & Description & Température & Densité & Ionisation \\
 & & [K] & [cm$^{-3}$] & \\
\midrule
Gaz moléculaire froid et dense (H$_2$) & Gaz froid et dense, composé principalement d'hydrogène moléculaire. & 10--20 & $>10^{2}$ & $<10^{-6}$ \\
Milieu neutre froid (CNM) & Gaz atomique froid, dense et principalement neutre. & 50--100 & 20--50 & $10^{-4}$ \\
Milieu neutre tiède (WNM) & Gaz atomique neutre, tiède et diffus. & 6000--10000 & 0.2--0.5 & $\approx 0.1$ \\
Milieu ionisé tiède (WIM) & Gaz tiède et diffus, partiellement ou entièrement ionisé. & $\sim 8000$ & 0.2--0.5 & 1.0 \\
Milieu ionisé chaud (HIM) & Gaz diffus, ionisé et de température élevée. & $\sim 10^{6}$ & $\sim 10^{-2}$ & 1.0 \\
\bottomrule
\end{tabularx}
\end{table}

\subsubsection{Les milieux neutres froid et tiède}

À l'équilibre thermique, lorsque le chauffage et le refroidissement se compensent, l'hydrogène atomique neutre peut se trouver dans deux phases thermiquement stables : le CNM et le WNM \citep{field1969, Wolfire2003}. Les nuages denses de CNM sont généralement entourés d'une enveloppe diffuse de WNM.\\\indent
Le gaz dont la température se situe entre celles de ces deux phases stables constitue le milieu neutre thermiquement instable (UNM), également désigné par LNM, pour \textit{lukewarm neutral medium}. Ce gaz peut se densifier et se refroidir pour rejoindre le CNM, ou au contraire se diluer et se réchauffer pour rejoindre le WNM \citep{field1969, Bialy2019}. Les observations ont confirmé la présence d'une quantité importante de gaz dans ce domaine thermiquement instable \citep{Heiles2003b, kanekar2003, Murray2018}.\\\indent
À sa densité habituelle, le WNM ne se transforme pas spontanément en CNM. Une compression isotherme d'un facteur deux à quatre suffit toutefois à déclencher une transition de phase. Le CNM se forme donc lorsque l'augmentation de pression fait entrer le WNM dans un régime thermiquement instable. Cette compression peut être provoquée par des phénomènes transitoires, tels que les explosions de supernovæ, les vents stellaires ou la turbulence à grande échelle \citep{Saury2014, Klessen2016}.

\begin{figure}[htbp]
\centering
\begin{minipage}{0.45\linewidth}
\centering
\includegraphics[width=.9\linewidth]{images/Monkey-Head-Nebula.pdf}
\subcaption{NGC 2174, ou nébuleuse de la Tête de singe}
\label{fig:monkeyhead}
\end{minipage}
\begin{minipage}{0.45\linewidth}
\centering
\includegraphics[width=.8\linewidth]{images/Part_of_the_Coalsack_Nebula_compresed.pdf}
\subcaption{Nébuleuse du Sac à charbon}
\label{fig:coalsack}
\end{minipage}
\caption{À gauche : NGC 2174, ou nébuleuse de la Tête de singe, illustrant la structure nuageuse du gaz H$_2$. Crédit : télescope spatial Hubble. À droite : la nébuleuse du Sac à charbon, un nuage moléculaire de la constellation de la Croix du Sud, visible à l'œil nu sous la forme d'une région sombre. Crédit : ESO.}
\end{figure}

\subsubsection{Le gaz moléculaire}

La phase moléculaire du MIS est dominée par le H$_2$. Elle contient aussi des molécules simples, comme le monoxyde de carbone (CO) et l'eau (H$_2$O), ainsi que des espèces organiques plus complexes, notamment le méthanol (CH$_3$OH) et le glycolaldéhyde (CH$_2$OHCHO), apparenté aux sucres \citep{sugar2000}. Les nuages moléculaires correspondent aux régions les plus denses du MIS, avec des densités en nombre généralement supérieures à $10^2$ cm$^{-3}$ (tableau~\ref{tab:ISMcomponents}), et peuvent contenir des grains de poussière. Ces densités restent néanmoins très inférieures à celle de l'atmosphère terrestre, de l'ordre de $10^{19}$ particules par centimètre cube. Les plus grandes structures moléculaires sont appelées nuages moléculaires géants (GMCs). Elles constituent les principaux sites de formation stellaire dans les galaxies \citep{Klessen2016}.

\subsubsection{Les milieux ionisés tiède et chaud}

Le milieu ionisé tiède (WIM) et le milieu ionisé chaud (HIM) constituent les phases ionisées du MIS. Le WIM présente une température d'environ 8000 K, entretenue par le rayonnement ionisant des étoiles massives. Le HIM atteint des températures de $\sim 10^6$ K et est principalement chauffé par les explosions de supernovæ \citep{Ferriere_2001}. Bien que ces phases jouent un rôle important dans le cycle de la matière interstellaire, cette étude porte surtout sur les phases atomique neutre et moléculaire, plus directement liées à la formation stellaire.

\subsection{Observer les phases de l'hydrogène atomique neutre}
\label{subsec:tracing_methods}

Les propriétés du MIS peuvent être étudiées à partir des raies d'émission et d'absorption de différentes espèces atomiques et moléculaires. Une raie d'émission apparaît lorsqu'un électron passe d'un niveau d'énergie élevé à un niveau plus bas et émet un photon dont la fréquence correspond à la différence d'énergie. Inversement, l'absorption d'un photon de fréquence appropriée permet à un électron d'atteindre un niveau supérieur. Les niveaux d'énergie étant quantifiés, chaque espèce possède des signatures spectrales caractéristiques.

\begin{figure}[t]
\centering
\begin{minipage}{0.9\linewidth}
\centering
\includegraphics[width=.9\linewidth]{images/Htransitions.pdf}
\end{minipage}
\caption{Diagramme des niveaux d'énergie de l'atome d'hydrogène. Le nombre quantique principal de l'électron est noté $n$. La différence d'énergie $\Delta E$ correspond à la transition de retournement du spin à l'origine de la raie à 21 cm. Crédit : DARA Foundation, Astrophysics -- Interstellar Gas, partie 2.}
\label{fig:Hseries}
\end{figure}

\subsubsection{La transition hyperfine à 21 cm}

L'hydrogène atomique neutre (\HI) est observable par sa transition hyperfine à 21 cm, située dans le domaine radio. L'état électronique fondamental de l'atome d'hydrogène comprend deux sous-niveaux d'énergie proches (figure~\ref{fig:Hseries}) : un état d'énergie plus élevée, où les spins de l'électron et du proton sont parallèles ($S=1$), et un état d'énergie plus faible, où ils sont antiparallèles ($S=0$). Le passage entre ces sous-niveaux, appelé transition de retournement du spin, correspond à une différence d'énergie $\Delta E = 5.87 \times 10^{-6}$ eV. Il produit un photon de fréquence $\nu = 1.4204$ GHz, soit une longueur d'onde $\lambda = 21.106$ cm \citep{ISF2011}.
Contrairement à la raie CO(1--0), optiquement épaisse et associée au gaz moléculaire dense où se forment les étoiles, l'émission à 21 cm est généralement optiquement mince dans la majeure partie du MIS. Elle permet de tracer le gaz atomique diffus loin des régions de formation stellaire, y compris au-delà du disque stellaire, où le gaz est trop diffus ou insuffisamment protégé pour abriter une phase moléculaire. Les relevés profonds à 21 cm montrent que les disques d'\HI\ s'étendent bien au-delà du disque optique. Le rayon $R_{\mathrm{HI}}$, défini à la densité de surface de $1\ \mathrm{M_\odot\,pc^{-2}}$, dépasse en moyenne le rayon optique $R_{25}$ d'un facteur $\sim 1.5$, avec des rapports parfois bien plus élevés \citep{Broeils1997, Wang2016}. Les observations pilotes récentes de WALLABY indiquent en outre qu'en moyenne environ deux tiers ($68\%$) de la masse totale d'\HI\ se trouvent à l'intérieur de $R_{25}$, le reste étant réparti à plus grand rayon dans un gaz de faible densité de surface \citep{Lee2025}. Ce réservoir étendu, peu actif en formation stellaire mais lié dynamiquement au disque, constitue un traceur des interactions, de l'accrétion et de la distribution de masse externe, difficilement accessibles par les seules observations optiques.

\subsubsection{Séparer les phases par l'absorption et l'émission}

La méthode classique pour déterminer les propriétés physiques de l'\HI\ et distinguer le CNM du WNM combine des observations en émission et en absorption à 21 cm \citep{Kulkarni1988, Dickey1990, Heiles2003a}. L'absorption est particulièrement sensible au CNM, dont la profondeur optique est élevée en raison de sa densité plus forte et de sa température plus basse que celles du WNM \citep{kanekar2003, Roy2013}.\\\indent
Dans le modèle à deux phases du MIS, la température de spin $T_s$, qui décrit les populations relatives des deux sous-niveaux hyperfins, est liée à la température cinétique $T_k$. Dans le CNM, les collisions sont assez fréquentes pour thermaliser la transition, de sorte que $T_s \approx T_k$, entre 40 et 200 K. Dans le WNM, la faible densité rend ce couplage collisionnel insuffisant : $T_s$ peut alors être sensiblement inférieure à $T_k$, à moins que le rayonnement Lyman $\alpha$ ne couple les deux températures \citep{Field1958, Liszt2001, Heiles2003b}.\\\indent
\citet{Heiles2003a, Heiles2003b} ont mené un relevé combinant absorption et émission devant 79 sources de continuum, le \textit{Millennium Arecibo 21 Centimeter Absorption-Line Survey}, et développé des méthodes de décomposition gaussienne pour séparer le CNM du WNM. Ils trouvent un maximum de la distribution des températures de spin du CNM vers 40 K et environ 60\% de l'\HI\ dans le WNM. Une fraction importante du WNM, au moins 48\%, se situe dans le domaine thermiquement instable, entre 500 et 5000 K.\\\indent
Les observations en absorption nécessitent toutefois des sources radio de continuum compactes et suffisamment brillantes en arrière-plan. Leur rareté limite fortement la couverture spatiale, surtout dans les galaxies extérieures \citep{kanekar2003, Murray2018}. Dans la Voie lactée, des relevés comme 21-SPONGE \citep{Murray2018} et THOR \citep{Beuther2016} ont augmenté le nombre de lignes de visée disponibles, mais ne permettent toujours pas de cartographier les deux phases dans l'ensemble des galaxies extérieures par la seule absorption.

\subsubsection{Décomposition gaussienne des spectres d'émission}
\label{subsubsec:gaussian_decomposition}

Pour contourner ces limites, d'autres méthodes cherchent à extraire l'information sur les phases à partir de la seule émission à 21 cm. Elles supposent généralement que le CNM et le WNM présentent des dispersions de vitesse caractéristiques différentes : les composantes du CNM sont étroites ($\sigma_v \sim 1$--6 km s$^{-1}$), en raison de leur basse température, tandis que celles du WNM sont plus larges ($\sigma_v \sim 6$--15 km s$^{-1}$), sous l'effet de l'agitation thermique et de la turbulence \citep{Young2003, deBlok2006, Ianjamasimanana2012}.\\\indent
\citet{Young2003} ont montré que ces composantes cinématiques pouvaient être identifiées dans les profils d'émission de sept galaxies naines proches. Leurs ajustements gaussiens distinguent des composantes étroites de $\sigma_v \sim 3$--5 km s$^{-1}$ et larges de $\sigma_v \sim 8$--13 km s$^{-1}$. Les premières sont regroupées dans des structures compactes souvent proches des régions de formation stellaire ; les secondes sont plus largement réparties. \citet{deBlok2006} ont ensuite ajusté des modèles à deux gaussiennes aux profils d'\HI\ de NGC 6822, en choisissant les modèles selon le $\chi^2$ réduit. Ils distinguent également des composantes étroites et larges, de dispersions moyennes respectives $\sim4$ et $\sim8$ km s$^{-1}$.\\\indent
L'ajustement de chaque profil reste cependant sensible au bruit, car les rapports signal sur bruit sont souvent faibles ou modérés. \citet{Ianjamasimanana2012} ont proposé la méthode des « superprofils » : les profils individuels de THINGS sont alignés en vitesse puis empilés pour obtenir un profil composite à haut rapport signal sur bruit. L'ajustement de deux gaussiennes révèle une composante étroite ($\sigma_v \sim 3.4$--8.6 km s$^{-1}$, moyenne $6.5\pm1.5$ km s$^{-1}$) et une composante large ($\sigma_v \sim 10.1$--24.3 km s$^{-1}$, moyenne $16.8\pm4.3$ km s$^{-1}$), interprétées comme des signatures du CNM et du WNM.\\\indent
Les travaux suivants ont étendu cette approche à l'étude des variations des deux composantes. \citet{Ianjamasimanana2015} ont empilé les profils d'\HI\ de 22 galaxies de THINGS, puis les ont décomposés en composantes gaussiennes étroite et large afin d'étudier la dispersion de vitesse et la densité de surface en fonction du rayon.\\\indent
D'autres développements visent à rendre la décomposition des profils individuels plus robuste et entièrement automatique. BAYGAUD (\textit{BAYesian GAUssian Decomposer}), proposé par \citet{Oh2019}, utilise une approche bayésienne par chaînes de Markov Monte-Carlo pour ajuster plusieurs modèles gaussiens à chaque profil et sélectionner le nombre optimal de composantes. Celles-ci peuvent ensuite être classées comme cinématiquement froides ou chaudes selon leur dispersion de vitesse. Par rapport à un ajustement classique fondé sur le $\chi^2$, BAYGAUD réduit la sensibilité aux paramètres initiaux et permet d'automatiser la décomposition et la classification pour de grands relevés d'\HI.\\\indent
\citet{Park2022} ont appliqué BAYGAUD à NGC 6822 et trouvé une association entre l'\HI\ cinématiquement froid et les régions de formation stellaire. \citet{KimOh2022} ont ensuite proposé de construire les superprofils en décomposant d'abord chaque spectre avec BAYGAUD, puis en alignant et en additionnant les composantes obtenues. Cette méthode vise à réduire les effets de l'asymétrie des profils, de la turbulence et du lissage par le faisceau, tout en conservant le gain en rapport signal sur bruit apporté par l'empilement.\\\indent
Plus récemment, \citet{Park2025} ont étudié les composantes d'\HI\ les plus étroites ($\sigma_v < 6$ km s$^{-1}$) comme traceurs du gaz associé à la formation stellaire dans sept galaxies proches de métallicités variées. Ces composantes présentent souvent une structure plus fragmentée que les composantes larges et une association morphologique avec les régions de formation stellaire. Cette association peut néanmoins s'atténuer aux meilleures résolutions spatiales, lorsque les décalages entre le gaz et les jeunes étoiles deviennent observables.

\subsubsection{Gaz cinématiquement froid et gaz thermiquement froid}
\label{subsubsec:kinematic_vs_thermal}

Il faut distinguer l'\HI\ « cinématiquement froid » de l'\HI\ « thermiquement froid ». Le CNM et le WNM sont définis par leurs températures physiques, que les observations combinées en absorption et en émission permettent de contraindre par la température de spin. La décomposition cinématique distingue, elle, des composantes selon leur dispersion de vitesse, qui comprend à la fois l'élargissement thermique, la turbulence et d'autres mouvements non thermiques.\\\indent
Une faible dispersion de vitesse est généralement associée au gaz froid, puisque l'élargissement thermique augmente avec la température, mais plusieurs effets compliquent cette interprétation. La turbulence peut élargir une raie intrinsèquement étroite ; le lissage par le faisceau produit un élargissement apparent comparable dans une galaxie en rotation \citep{Ianjamasimanana2012, Caldu-Primo2013, Stilp2013}. De plus, plusieurs nuages froids à des vitesses différentes sur une même ligne de visée peuvent produire un profil composite plus large que chacune de ses composantes.\\\indent
Nous utilisons donc les termes « cinématiquement froid » et « cinématiquement chaud » pour les composantes issues de la décomposition gaussienne des spectres d'émission, sans les assimiler exactement au CNM et au WNM thermiques. Cette distinction est importante pour interpréter leurs liens avec la formation stellaire : le lien entre les raies étroites d'\HI\ et les conditions de formation des nuages moléculaires dépend notamment de l'instabilité thermique, de la protection par la poussière et de l'effondrement gravitationnel \citep{Schaye2004, Krumholtz2011, Klessen2016}.

\subsubsection{Tracer l'hydrogène moléculaire}

Contrairement à l'\HI, le H$_2$ est très difficile à observer directement aux températures habituelles des nuages moléculaires ($\sim$10--40 K). Cette molécule diatomique homonucléaire ne possède pas de moment dipolaire électrique permanent et son moment d'inertie est faible. Ses premières transitions rotationnelles nécessitent des températures d'excitation proches de 540 K, bien supérieures à celles des nuages moléculaires \citep{Ferriere_2001, ISF2011}.\\\indent
Le H$_2$ peut être détecté par ses raies d'absorption ultraviolettes, mais cette méthode nécessite une étoile chaude et suffisamment brillante sur la même ligne de visée, ce qui empêche une cartographie étendue. Le monoxyde de carbone (CO), relativement abondant et dont les transitions rotationnelles sont excitées à basse température, sert donc couramment de traceur du gaz moléculaire. La densité de colonne de H$_2$ est déduite du CO à l'aide d'un facteur de conversion, noté $X_{\mathrm{CO}}$ ou $\alpha_{\mathrm{CO}}$, qui dépend notamment de la métallicité \citep{Bolatto2013}.\\\indent
Cette méthode présente des limites importantes à faible métallicité. Plus facilement photodissocié que le H$_2$, le CO ne subsiste que dans les régions les plus denses et les mieux protégées des nuages. Une fraction parfois importante du H$_2$ échappe alors aux observations du CO : on parle de gaz moléculaire « sombre en CO » \citep{Bolatto2013, Wolfire2010}. Ce phénomène est particulièrement marqué dans les enveloppes des nuages moléculaires et dans les galaxies pauvres en métaux, où la poussière protège moins efficacement le gaz.

\section{La formation stellaire et le milieu interstellaire}
\label{sec:SF_and_ism}
\subsection{Vue d'ensemble de la formation stellaire}

Les étoiles naissent dans les nuages moléculaires lorsque la gravité l'emporte sur le soutien assuré par la pression interne. La longueur de Jeans ($r_J$) caractérise cette condition et représente le rayon minimal à partir duquel un nuage sphérique de densité uniforme devient gravitationnellement instable :
\begin{equation}
    r_J \sim \frac{a_0}{(G\rho_0)^{1/2}}
\end{equation}
\noindent où $G$ est la constante gravitationnelle, $\rho_0$ la densité initiale du gaz et $a_0$ la vitesse du son. Cette longueur étant inversement proportionnelle à la racine carrée de la densité, une compression externe peut déclencher l'effondrement : explosion de supernova voisine, onde de choc, collision entre nuages ou interaction entre galaxies. L'effondrement peut aussi commencer lorsqu'un nuage dépasse la masse de Jeans \citep{Jeans1902, ISF2011}.\\\indent
Au cours de la contraction, certaines parties du nuage deviennent à leur tour instables et se fragmentent en structures plus petites et plus denses. Chaque fragment continue à se contracter jusqu'à ce qu'il devienne assez opaque pour piéger le rayonnement émis lors du refroidissement. Un cœur chaud, soutenu par la pression, se forme alors et poursuit son évolution jusqu'à la formation d'une protoétoile \citep{ISF2011, Klessen2016}.

\subsection{La transition \HI--H$_2$ et son rôle dans la formation stellaire}

Le MIS neutre diffus étant surtout atomique, alors que les nuages moléculaires géants où se forment les étoiles sont principalement constitués de H$_2$, leur évolution implique une étape de conversion de l'hydrogène atomique en hydrogène moléculaire. Cette \textbf{transition \HI--H$_2$} joue un rôle fondamental dans le cycle de la matière interstellaire et la formation stellaire \citep{Lee2012, Sternberg2014, Klessen2016}.\\\indent
Le taux de formation stellaire est étroitement lié à cette transition \citep{Krumholtz2011}, car les conditions favorables à la formation du H$_2$ sont proches de celles requises pour former des étoiles. Le gaz doit se refroidir suffisamment pour que la gravité surmonte la pression thermique. Le champ de rayonnement interstellaire (ISRF) peut toutefois réchauffer le gaz et s'opposer à ce refroidissement. La présence de H$_2$ indique souvent celle de grains de poussière, qui jouent un double rôle : ils catalysent la formation du H$_2$ à leur surface et absorbent le rayonnement interstellaire, limitant le chauffage du gaz tout en protégeant les molécules contre la photodissociation \citep{Klessen2016}.\\\indent
Dans les environnements riches en métaux, le refroidissement est dominé par les raies des éléments lourds. Le H$_2$ peut néanmoins devenir un agent de refroidissement à faible métallicité ($< 0.1\,Z_\odot$) ou lorsque le champ de rayonnement est faible, permettant la formation stellaire dans des nuages primordiaux ou pauvres en métaux \citep{Glover2014}.

\subsection{Les lois empiriques de formation stellaire}
\label{sec:sfr_ism_relationship}

Comme nous l'avons vu, les différentes phases du MIS fournissent le gaz nécessaire à la formation des étoiles. Comprendre leurs liens est essentiel pour modéliser l'évolution des galaxies. À l'échelle galactique, la complexité des processus en jeu limite cependant les prédictions purement théoriques. Les corrélations observées entre les propriétés du gaz et la formation stellaire sont donc souvent décrites par des relations empiriques en loi de puissance, appelées lois de formation stellaire.\\\indent
\citet{Schmidt1959, Schmidt1963} ont proposé de relier la densité volumique du taux de formation stellaire $f$ à la densité volumique du gaz $\rho_{\mathrm{gas}}$ :
\begin{equation}
    f \propto \rho_{\mathrm{gas}}^n
\end{equation}
\noindent Une formulation très utilisée relie ensuite le taux de formation stellaire par unité de surface, $\Sigma_{\mathrm{SFR}}$, à la densité de surface du gaz, $\Sigma_{\mathrm{gas}}$ :
\begin{equation}
    \Sigma_{\mathrm{SFR}} = A\,(\Sigma_{\mathrm{gas}})^N
    \label{eq:SFLaw}
\end{equation}
\noindent Cette relation est appelée loi de Kennicutt--Schmidt. L'indice $N$ est l'exposant de la loi de puissance, ou sa pente en coordonnées logarithmiques, et $A$ son coefficient de normalisation. Les densités de surface ont l'avantage d'être directement accessibles dans les galaxies extérieures, contrairement aux densités volumiques, qui nécessitent des hypothèses sur la distribution du gaz le long de la ligne de visée. On exprime généralement $\Sigma_{\mathrm{SFR}}$ en $M_\odot\,\mathrm{kpc}^{-2}\,\mathrm{yr}^{-1}$ et $\Sigma_{\mathrm{gas}}$ en $M_\odot\,\mathrm{pc}^{-2}$.\\\indent
La valeur de $N$ renseigne sur la régulation de la formation stellaire. Pour $N = 1$, le taux de formation stellaire augmente proportionnellement à la quantité de gaz. Pour $N > 1$, les régions les plus denses forment les étoiles plus efficacement ; pour $N < 1$, l'efficacité diminue lorsque la densité de gaz augmente.\\\indent
\citet{Kennicutt1998} a étudié la relation entre $\Sigma_{\mathrm{SFR}}$ et $\Sigma_{\textrm{H\textsc{i}}+\mathrm{H}_2}$, moyennées sur le disque, pour 61 galaxies spirales normales et 36 galaxies à flambée de formation stellaire. Il trouve $N = 2.47 \pm 0.39$ pour les seules spirales normales et $N = 1.4 \pm 0.15$ pour l'échantillon complet. \citet{delosReyes2019} ont ensuite révisé cette dernière valeur à $N = 1.41 \pm 0.07$. L'indice $N \approx 1.4$ est largement utilisé dans les simulations de formation stellaire, d'évolution galactique et de cosmologie \citep{SchayetoDalla2008, FOREVER22, FLAMINGO}.\\\indent
Cet indice dépend de l'échelle spatiale, du type de galaxie et de la composante gazeuse étudiée. \citet{Boissier2003} trouvent $N \approx 2$ dans 16 galaxies. \citet{Misiriotis2006} obtiennent $N = 2.18 \pm 0.20$ pour la Voie lactée, et \citet{Heyer2004}, $N = 3.3 \pm 0.1$ pour M33. Pour le gaz moléculaire seul, \citet{Leroy2005} trouvent $N = 1.3 \pm 0.1$ dans les galaxies naines, tandis que \citet{Fbigiel2008} obtiennent $N = 1.0 \pm 0.2$ pour $\Sigma_{\mathrm{H}_2}$ dans les spirales, ce qui suggère une relation presque linéaire entre le gaz moléculaire et la formation stellaire.\\\indent
Une autre manière de décrire ce lien est l'efficacité de formation stellaire (SFE) :
\begin{equation}
    \mathrm{SFE} = \frac{\Sigma_{\mathrm{SFR}}}{\Sigma_{\mathrm{gas}}}
    \label{eq:sfe}
\end{equation}
\noindent Exprimée en yr$^{-1}$, elle représente la fraction du gaz convertie en étoiles par unité de temps. Son inverse est le temps d'épuisement du gaz ($t_{\mathrm{dep}} = 1/\mathrm{SFE}$), c'est-à-dire la durée nécessaire pour consommer le réservoir si la formation stellaire se maintient au taux actuel \citep{Leroy2008, Yu-Hsuan2024}.

\subsection{Les seuils de formation stellaire et le rôle du gaz froid}

La présence de gaz ne suffit pas à garantir la formation d'étoiles. Celle-ci augmente de façon monotone aux fortes densités de gaz, mais diminue nettement aux faibles densités. On définit ainsi une densité de surface critique, $\Sigma_{\mathrm{crit}}$, en dessous de laquelle le gaz forme peu efficacement des étoiles. Les seuils observés se situent généralement entre 5 et 10 $M_\odot\,\mathrm{pc}^{-2}$ \citep{Kennicutt1998, Martin2001}.\\\indent
Plusieurs mécanismes ont été proposés pour expliquer ces seuils : instabilités gravitationnelles à grande échelle, cisaillement, métallicité ou turbulence \citep{Leroy2008, Sternberg2014, Federrath2015}. \citet{Schaye2004} montre en particulier que la diminution de la dispersion thermique lors de la transition du WNM vers le CNM peut déclencher une instabilité gravitationnelle et conduire à l'effondrement du gaz. Il propose d'utiliser l'apparition d'\HI\ froid, accompagnée d'une instabilité thermogravitationnelle, comme critère de formation stellaire, avec $\Sigma_{\mathrm{crit}} \approx 3$--10 $M_\odot\,\mathrm{pc}^{-2}$.\\\indent
Les travaux antérieurs supposant que le CNM se trouve là où la formation du H$_2$ est efficace, nous cherchons à tester cette hypothèse en examinant les relations entre la formation stellaire et les différentes phases du gaz, en particulier l'\HI\ cinématiquement froid. La relation de Kennicutt--Schmidt est appliquée aux différentes composantes gazeuses. Le seuil fondé sur la formation de gaz froid proposé par \citet{Schaye2004} est également examiné, et les efficacités de formation stellaire de l'\HI, du H$_2$ et de l'\HI\ cinématiquement froid sont comparées aux études antérieures.
